\documentclass[12pt]{article}
\usepackage{amsfonts,amsthm,amsmath,amssymb,mathtools}
\usepackage{mathrsfs}
\usepackage{enumitem}
\usepackage{tikz-cd}
\usepackage{geometry}
\usepackage{graphicx}
\IfFileExists{mlmodern.sty}{
    \usepackage{mlmodern}
}{}

\geometry{letterpaper, portrait, margin=1in}

\setcounter{section}{0}

\renewcommand{\thesection}{\S\arabic{section}}
\renewcommand{\thesubsection}{\arabic{section}}
\DeclareMathOperator{\im}{im}
\DeclareMathOperator{\id}{id}

\newtheorem{thm}{Theorem}
\newenvironment{theorem}[1]{
	\renewcommand{\thethm}{#1}
	\begin{thm}
		}{\end{thm}}

\newtheorem{lma}{Lemma}
\newenvironment{lemma}[1]{
	\renewcommand{\thelma}{#1}
	\begin{lma}
		}{\end{lma}}

\theoremstyle{definition}
\newtheorem{ex}{Problem}
\newenvironment{exercise}[1]{
	\renewcommand{\theex}{#1}
	\begin{ex}
		}{\end{ex}}

\title{Functional Analysis --- Homework 7}
\author{Connor Mooneyhan}
\date{}

\begin{document}

\maketitle

\begin{ex}
	Show that $C^0[0, 1]$ in the supremum norm is not an inner product space (that is, the norm cannot come from an inner product).
\end{ex}
\begin{proof}

\end{proof}

\begin{ex}
	In this problem we review the Gram--Schmidt process: Let $x_1, x_2, \dots$ be a sequence of linearly independent vectors in an inner product space.
	Define vectors inductively by: \begin{align*}
		 & e_1 = x_1 / \|x_1\|;                                                                                                \\
		 & f_n = x_n - \sum_{j=1}^{n-1} \langle x_n, e_j \rangle e_j, \hspace{1.5em} \text{and} \hspace{1.5em} e_n = f_n / \|f_n\|, \hspace{1.5em} \text{for } n \geq 2.
	\end{align*}
	Show that $\lbrace e_n \rbrace$ is an orthonormal sequence with the property that the span of $\lbrace x_1, x_2, \dots, x_n \rbrace$ is the same as the span of $\lbrace e_1, e_2, \dots, e_n \rbrace$ for each $n$.
\end{ex}
\begin{proof}

\end{proof}

\begin{ex}
	Apply the Gram--Schmidt process (Problem 2) to the three vectors $\lbrace 1, x, x^2 \rbrace$ in the space $L^2[-1, 1]$.
	Use your answer to find the distance from $x^3$ to the span of $\lbrace 1, x, x^2 \rbrace$, i.e., \begin{equation*}
		\min_{a, b, c \in \mathbb{C}} \int_{-1}^1 \left| x^3 - a - bx - cx^2 \right|^2 dx.
	\end{equation*}
	Remark: when the Gram--Schmidt process is applied to the sequence $1, x, x^2, x^3, \dots$, the resulting vectors are called the \textit{Legendre polynomials}.
\end{ex}
\begin{proof}

\end{proof}

\begin{ex}
	Complete the argument outlined below, originally due to Riesz, to give a proof of the completeness of $L^2[a, b]$.
	Let $\lbrace f_n \rbrace$ be a Cauchy sequence in $L^2[a, b]$.
	\begin{enumerate}[label=(\alph*)]
		\item Show that $\lbrace \|f_n\| \rbrace$ is a Cauchy sequence of scalars, and hence a bounded sequence, say $\|f_n\| \leq M$.
		\item Show that for fixed $g \in L^2[a, b]$, the sequence $\int_a^b f_ng \, dx$ is Cauchy and its limit, call it $\alpha_g$, satisfies the bound $|\alpha_g| \leq M\|g\|$.
		\item Define $\Lambda : L^2[a, b] \to \mathbb{C}$ by $\Lambda(g) = \alpha_g$.
		      Use the Riesz representation theorem to show that there exists $F \in L^2[a, b]$ with $\alpha_g = \int_a^b Fg \, dx$ for all $g \in L^2[a, b]$.
		\item Show that $f_n$ converges to $F$ in $L^2[a, b]$.
	\end{enumerate}
\end{ex}
\begin{proof}\
	\begin{enumerate}[label=(\alph*)]
		\item
		\item
		\item
		\item
	\end{enumerate}
\end{proof}

\begin{ex}
	A sequence $\lbrace h_n \rbrace$ in a Hilbert space $\mathscr{H}$ is said to \textit{converge weakly} to $h \in \mathscr{H}$ if \begin{equation*}
		\lim_{n \to \infty} \langle h_n, g \rangle = \langle h, g \rangle
	\end{equation*}
	for every $g \in \mathscr{H}$.
	\begin{enumerate}[label=(\alph*)]
		\item If $\lbrace e_n \rbrace$ is an orthonormal sequence in $\mathscr{H}$, show that $e_n \to 0$ weakly.
		\item Show that if $h_n \to h$ in norm, then $h_n \to h$ weakly.
		      Show that the converse is false, but that if $h_n \to h$ weakly and $\|h_n\| \to \|h\|$, then $h_n \to h$ in norm.
	\end{enumerate}
\end{ex}
\begin{proof}\
	\begin{enumerate}[label=(\alph*)]
		\item
		\item
	\end{enumerate}
\end{proof}

\end{document}
